% CC BY 4.0 License
% Copyright (c) 2026 Nanum Learn

\documentclass[a4paper,11pt]{article}

\usepackage[T1]{fontenc}
\usepackage[utf8]{inputenc}
\usepackage[english]{babel}

\usepackage{amsmath}
\usepackage{amssymb}
\usepackage{amsthm}
\usepackage{mathtools}

\usepackage[margin=1in]{geometry}
\usepackage{lastpage}
\usepackage{fancyhdr}
\fancyhead[L]{NL01CHEM01 Session 2 Problem Set}
\fancyhead[R]{Nanum Learn}
\fancyfoot[C]{}
\pagestyle{fancy}
\usepackage[skip=10pt]{parskip}
\setlength{\headheight}{13.59999pt}
\setlength{\headsep}{8pt}

\usepackage{titlesec}
\titleformat{\section}
    {\normalfont\LARGE\bfseries\color{BrickRed}}
    {\S\! \thesection}
    {0.5em}
    {}
\titleformat{\subsection}
    {\normalfont\Large\bfseries}
    {\thesubsection}
    {0.5em}
    {}

\usepackage{enumitem}
\usepackage[dvipsnames]{xcolor}


\newcommand{\easy}{(\ensuremath{\bullet \circ \circ})}
\newcommand{\medium}{(\ensuremath{\bullet \bullet \circ})}
\newcommand{\hard}{(\ensuremath{\bullet \bullet \bullet})}

\newtheoremstyle{listprobstyle}
    {\parsep}{\parsep}
    {\normalfont}
    {}
    {\bfseries}
    {}
    { }
    {{\thmname{#1}~\thmnumber{#2}.}~
        \thmnote{\textcolor{BrickRed}{\bfseries #3}}\\}
\theoremstyle{listprobstyle}
\newtheorem{problem}{Problem}



\title{\textbf{NL01CHEM01 Session 2 Problem Set}}
\date{10 November 2026}
\author{William Kang}

\begin{document}

\maketitle
\begin{center}
    \small\textit{Edited and Formatted by Enoch Yu}
\end{center}

\vspace{\baselineskip}
\begin{center}
\begin{minipage}{0.7\textwidth}
These are six chemistry practice problems written by Chemistry
director, William Kang. The problems are arranged roughly in
increasing difficulty with subjective ratings, and covers the topics
from William's presentation today. Enjoy!
\end{minipage}
\end{center}
\vspace{\baselineskip}

\section{Problems}
\begin{problem}[Outer Electrons from Column \easy]
How many outer electrons do elements in each column have?
\vspace{-0.6\baselineskip}
\begin{enumerate}[label=\alph*),itemsep=0pt,parsep=0pt,topsep=0pt]
\item Column 1 (lithium family): \underline{\phantom{XXX}}
\item Column 14 (carbon family): \underline{\phantom{XXX}}
\item Column 17 (chlorine family): \underline{\phantom{XXX}}
\end{enumerate}
\end{problem}

\begin{problem}[Column from Outer Count \easy]
Identify the periodic table column from the number of outer electrons.
\vspace{-0.6\baselineskip}
\begin{enumerate}[label=\alph*),itemsep=0pt,parsep=0pt,topsep=0pt]
\item 2 outer electrons: Column \underline{\phantom{XXX}}
\item 6 outer electrons: Column \underline{\phantom{XXX}}
\item 7 outer electrons: Column \underline{\phantom{XXX}}
\end{enumerate}
\end{problem}

\begin{problem}[Write the Ion \easy]
For each atom, write the ion it typically forms.
\vspace{-0.6\baselineskip}
\begin{enumerate}[label=\alph*),itemsep=0pt,parsep=0pt,topsep=0pt]
\item Sodium: \underline{\phantom{XXXXXX}}
\item Magnesium: \underline{\phantom{XXXXXX}}
\item Chlorine: \underline{\phantom{XXXXXX}}
\item Oxygen: \underline{\phantom{XXXXXX}}
\end{enumerate}
\end{problem}

\begin{problem}[Ion or Share? \easy]
Label each compound I (ionic) or C (covalent).
\vspace{-0.6\baselineskip}
\begin{enumerate}[label=\alph*),itemsep=0pt,parsep=0pt,topsep=0pt]
\item Salt ($\mathrm{NaCl}$): (\phantom{X})
\item Water ($\mathrm{H_2O}$): (\phantom{X})
\item $\mathrm{MgCl_2}$: (\phantom{X})
\item $\mathrm{CO_2}$: (\phantom{X})
\end{enumerate}
\end{problem}

\newpage
\begin{problem}[Write the Formula \easy]
Write the neutral chemical formula formed by each pair of ions.
\vspace{-0.6\baselineskip}
\begin{enumerate}[label=\alph*),itemsep=0pt,parsep=0pt,topsep=0pt]
\item Sodium ($\mathrm{Na^+}$) and chlorine ($\mathrm{Cl^-}$):
    \underline{\phantom{XXXXXX}}
\item Magnesium ($\mathrm{Mg^{2+}}$) and oxygen ($\mathrm{O^{2-}}$):
    \underline{\phantom{XXXXXX}}
\end{enumerate}
\end{problem}

\begin{problem}[Count Atoms in a Formula \easy]
Consider the chemical formula $\mathrm{Al_2O_3}$.
\vspace{-0.6\baselineskip}
\begin{enumerate}[label=\alph*),itemsep=0pt,parsep=0pt,topsep=0pt]
\item How many aluminum atoms? \underline{\phantom{XXX}}
\item How many oxygen atoms? \underline{\phantom{XXX}}
\item What is the Al:O ratio? \underline{\phantom{XXXX}}
\end{enumerate}
\end{problem}

\begin{problem}[Lose or Gain? \easy]
For each ion, state whether the atom lost or gained electrons.
\vspace{-0.6\baselineskip}
\begin{enumerate}[label=\alph*),itemsep=0pt,parsep=0pt,topsep=0pt]
\item $\mathrm{Na^+}$: \underline{\phantom{XXXXXX}}
\item $\mathrm{Cl^-}$: \underline{\phantom{XXXXXX}}
\item $\mathrm{Mg^{2+}}$: \underline{\phantom{XXXXXX}}
\item $\mathrm{O^{2-}}$: \underline{\phantom{XXXXXX}}
\end{enumerate}
\end{problem}

\begin{problem}[True or False? \easy]
Write T/F next to each statement.
\vspace{-0.6\baselineskip}
\begin{enumerate}[label=\alph*),itemsep=0pt,parsep=0pt,topsep=0pt]
\item $\mathrm{MgCl_2}$ is neutral overall. (\phantom{X})
\item $\mathrm{Na^+}$ has more electrons than protons. (\phantom{X})
\item Subscripts in chemical formulas indicate charge. (\phantom{X})
\item Elements in the same column have the same number of electron
    shells. (\phantom{X})
\end{enumerate}
\end{problem}

\begin{problem}[A Trickier Pair \medium]
Calcium forms a $2+$ ion ($\mathrm{Ca^{2+}}$), and hydroxide is a
$1-$ ion ($\mathrm{OH^-}$).
\vspace{-0.6\baselineskip}
\begin{enumerate}[label=\alph*),itemsep=0pt,parsep=0pt,topsep=0pt]
\item Write the neutral chemical formula:
    \underline{\phantom{XXXXXXX}}
\item How many calcium ions are needed?
    \underline{\phantom{XXX}}
\item How many hydroxide ions are needed?
    \underline{\phantom{XXX}}
\end{enumerate}
\end{problem}

\begin{problem}[Water in Lewis Form \medium]
Consider the Lewis structure of water ($\mathrm{H_2O}$).
\vspace{-0.6\baselineskip}
\begin{enumerate}[label=\alph*),itemsep=0pt,parsep=0pt,topsep=0pt]
\item How many shared electron pairs are there?
    \underline{\phantom{XXXX}}
\item How many lone pairs are on the oxygen atom?
    \underline{\phantom{XXXX}}
\item How many electrons surround the oxygen atom?
    \underline{\phantom{XXXX}}
\end{enumerate}
\end{problem}
\end{document}


